% =============================================================================
% TEMPLATE-ONLY RESEARCH WRITING GUIDE - ENGLISH
% Disable it in configuration.tex with:
% \providecommand{\ISISResearchWritingGuideChoice}{disabled}
% This chapter adapts general research-paper advice to a longer thesis. It is
% teaching material for the template and should not appear in the final thesis.
% =============================================================================

\chapter{Template Guide: Writing a Strong Research Thesis}\label{ch:writing-guide-en}

\begin{keypoint}[Template-only chapter]
This chapter demonstrates how to turn a research idea into a readable thesis.
It is not part of the example research project. Disable it in
\texttt{configuration.tex} before submission by setting
\texttt{\textbackslash ISISResearchWritingGuideChoice} to \texttt{disabled}.
\end{keypoint}

The guidance is a thesis-oriented adaptation of Simon Peyton Jones's Microsoft
Research presentation \emph{How to Write a Great Research Paper}. The source
presentation is available from
\href{https://www.microsoft.com/en-us/research/wp-content/uploads/2016/07/How-to-write-a-great-research-paper.pdf}{Microsoft Research}.
A thesis is longer than a conference paper and may contain several studies, but
it still benefits from one central claim, an explicit narrative, refutable
contributions, concrete examples, and evidence linked to every important claim.

\section{Start from one central thesis claim}
A strong thesis can usually be summarised as one sentence that states what was
learned, built, demonstrated, or explained. Several contributions may support
that sentence, but they should not feel like unrelated projects.

\begin{researchquestion}[Central thesis claim -- example]
A modular edge-AI design can reduce response time and resource use while
preserving useful predictive performance, provided that sensing, inference,
security policy, and presentation are evaluated as separate components.
\end{researchquestion}

\begin{keypoint}[Adaptation from paper to thesis]
For a paper, a single sharp idea may be enough. For a thesis, keep one central
claim and organise two to four subordinate contributions around it. Each
contribution should be specific enough that a reader could test whether the
thesis actually delivered it.
\end{keypoint}

A useful narrative for a technical thesis is:
\begin{enumerate}
  \item a concrete problem that the reader can picture;
  \item why the problem matters and which requirement is unmet;
  \item the thesis idea and the assumptions under which it applies;
  \item the method or system that realises the idea;
  \item evidence that supports or limits the claim;
  \item a fair comparison with previous approaches;
  \item conclusions, limitations, and next questions.
\end{enumerate}

\section{Example page: a four-sentence abstract}
The abstract should let a reader recover the problem, the approach, the main
evidence, and the contribution without opening the rest of the thesis. The
following example uses four sentences; the numbers are placeholders.

\begin{frontmatterbox}
\textbf{Problem.} Edge devices used in smart environments must respond quickly
while operating with limited memory and a clearly defined security boundary.
\textbf{Approach.} This thesis designs a modular prototype that separates
sensing, local inference, policy enforcement, and presentation, and evaluates
each component under a reproducible protocol. \textbf{Evidence.} In the example
evaluation, the optimised configuration reaches 91.0\% accuracy with a median
latency of 29~ms and 241~MB of memory, compared with 72.0\%, 84~ms, and 410~MB
for the baseline. \textbf{Contribution.} The results show how an evidence-linked
architecture can make performance, resource, and security trade-offs explicit;
all values must be replaced with the student's real findings.
\end{frontmatterbox}

\begin{keypoint}[Abstract test]
After reading only the abstract, a supervisor should be able to underline one
problem, one method, one principal result, and one contribution. Remove generic
claims such as ``technology is increasingly important'' unless they directly
establish the research problem.
\end{keypoint}

\section{Example page: introduce the problem with a running case}
Start with a small case rather than a broad survey. The case gives technical
terms a purpose and helps the reader remember the central tension.

\begin{resultbox}[Example opening paragraph]
At 02:14, a camera in an industrial room detects motion near a restricted
cabinet. The network link is temporarily unavailable, so a cloud-only system
cannot return a decision within the 50~ms response budget. A local model can
respond in time, but only if its memory use remains below 300~MB and its false
alarm rate stays within the operational limit. This thesis asks whether a
modular edge pipeline can satisfy those three constraints together.
\end{resultbox}

The paragraph does four jobs: it presents a concrete situation, exposes the
constraint, identifies the gap, and leads naturally to a research question. It
does not begin with several pages of general statements about the importance of
artificial intelligence.

\section{Example page: state refutable contributions}
Write the contribution list early, then use it to decide what belongs in the
thesis. Avoid contributions that merely say that a system is described or that
an issue is studied. Prefer claims that identify an artefact, property,
comparison, or measured outcome.

\begin{keypoint}[Example contribution list]
\begin{enumerate}
  \item We define a measurable requirement that combines latency, memory, and
        security constraints, and we document its assumptions in the problem
        formulation chapter.
  \item We implement a four-stage architecture whose components can be tested
        independently, and we provide the interface and configuration details
        needed to rebuild it.
  \item We compare the prototype with a stated baseline over repeated trials
        and report accuracy, latency, memory, uncertainty, and failure cases.
  \item We identify the operating conditions in which the approach does not
        meet the target, rather than presenting only the best configuration.
\end{enumerate}
\end{keypoint}

\subsection{Link every claim to evidence}
The introduction makes promises; the body must provide evidence for each one.
A claim-evidence matrix is a simple planning device.

\begin{table}[htbp]
\centering
\caption{Example claim-evidence map for a thesis.}
\label{tab:guide-claim-evidence-en}
\begin{tabularx}{\textwidth}{p{0.23\textwidth}Yp{0.27\textwidth}}
\toprule
Claim & Required evidence & Planned location\\
\midrule
The prototype meets the latency target & Repeated measurements, distribution,
comparison with the baseline & Evaluation protocol, latency figure, and raw
results in the appendix\\
The architecture is modular & Interface definition, component tests, and a
replacement experiment & Design chapter and component-level tests\\
The method is robust to one missing input & Predefined corruption conditions,
performance loss, and uncertainty & Robustness experiment and failure analysis\\
The result is reproducible & Versioned data, code, parameters, seeds, hardware,
and a reproduction command & Method chapter, appendix, and optional checklist\\
\bottomrule
\end{tabularx}
\end{table}

Use forward references from each contribution to the section, table, theorem,
case study, or experiment that supports it. A mechanical paragraph listing the
order of chapters is less useful than a narrative that connects claims to
evidence.

\section{Example page: intuition before formal detail}
Explain the idea as if drawing it on a whiteboard. Introduce notation only after
the reader understands what the components do and why they are needed.

\begin{figure}[htbp]
\centering
\pipelinebox{Observe\\quality}
\pipelinearrow
\pipelinebox{Estimate\\reliability}
\pipelinearrow
\pipelinebox{Weight\\evidence}
\pipelinearrow
\pipelinebox{Predict and\\report limits}
\caption{A whiteboard-level explanation before equations or implementation details.}
\label{fig:guide-intuition-en}
\end{figure}

\begin{keypoint}[Before and after]
\textbf{Too early:} ``Let $z_m=f_m(x_m)$ and optimise a regularised multimodal
objective.'' The notation arrives before the reader knows the problem.\par
\textbf{Better:} ``Each sensor proposes an answer and reports how trustworthy
its current signal is. The fusion module gives less influence to a sensor that
is noisy or absent. We formalise that behaviour after the running example.''
\end{keypoint}

Once the intuition is stable, define symbols, assumptions, algorithms, and
proofs precisely. Examples do not replace rigour; they prepare the reader to
understand it.

\section{Use active, direct sentences}
Prefer a sentence that names the actor and action. For example, write ``We ran
30 independent trials'' rather than ``Thirty trials were run'', and write ``The
optimised model used 241~MB'' rather than ``A reduction in memory utilisation
was observed''. Use the passive voice when the actor is genuinely irrelevant,
not as the default style.

\clearpage
\section{Example page: place related work where it helps}
A thesis often needs more background than a paper, and a degree programme may
require an early literature review. Keep only the concepts needed to understand
the problem near the beginning. Place the dense comparison with alternatives
after the reader understands the proposed idea, or revisit the literature in a
later positioning chapter.

\begin{table}[htbp]
\centering
\caption{Example synthesis table: compare assumptions rather than listing papers.}
\label{tab:guide-related-en}
\begin{tabularx}{\textwidth}{Ycccc}
\toprule
Approach & Local decision & Missing input & Calibration & Resource report\\
\midrule
Cloud-only baseline & No & Service failure & Optional & Server-side only\\
Early fusion & Yes & Limited & Rare & Sometimes\\
Late fusion & Yes & Moderate & Optional & Often\\
Thesis approach & Yes & Explicit & Required & Device and server\\
\bottomrule
\end{tabularx}
\end{table}

Describe previous work generously and accurately. State what an approach does
well, the assumptions under which it succeeds, and the precise dimension on
which the thesis extends or differs from it. Acknowledging weaknesses in the
new approach makes the comparison more credible.

\section{Example page: report evidence, not only winning numbers}
A results chapter should separate observation from interpretation.

\begin{resultbox}[Observation]
Across five illustrative seeds, median latency decreased from 84~ms to 29~ms.
The interquartile range decreased from 12~ms to 5~ms. One hardware configuration
exceeded the target when thermal throttling began after 20 minutes.
\end{resultbox}

\begin{keypoint}[Interpretation]
The lower median and narrower spread support the responsiveness claim on the
tested hardware. The throttling failure limits the claim to workloads and
cooling conditions represented by the protocol. The thesis should not infer
performance on untested devices.
\end{keypoint}

Report uncertainty, negative results, protocol deviations, and alternative
explanations. Distinguish exploratory findings from tests that were planned in
advance. A limitation is not an apology; it defines the boundary of the result.

\clearpage
\section{Revision page: test the reader's route through the thesis}
A reader should be able to follow the central claim by reading the title,
abstract, introduction, chapter openings, figure captions, result summaries,
and conclusions. Use the following review pass before polishing typography.

\begin{enumerate}
  \item Write the central thesis claim in one sentence.
  \item Check that the abstract states problem, approach, result, and contribution.
  \item Replace a broad opening with one concrete running example.
  \item Make every contribution refutable and point it to evidence.
  \item Explain the intuition before notation, algorithms, or implementation detail.
  \item Move dense related-work comparison until the reader knows the thesis idea,
        unless institutional rules require a different order.
  \item Separate measured observations from interpretation and speculation.
  \item Report uncertainty, limitations, negative results, and scope conditions.
  \item Prefer active, direct language and informative figures.
  \item Ask both an expert and a non-expert reader where they became lost, then
        revise the text rather than blaming the reader.
\end{enumerate}

\begin{keypoint}[Thesis-specific final check]
The Microsoft Research presentation addresses papers, not university
regulations. Preserve required chapters, declarations, formatting, ethics
statements, and assessment criteria from the degree programme. Use the narrative
principles here inside those institutional constraints.
\end{keypoint}
